\documentclass{techbrief}

\providecommand{\Dom}{\operatorname{Dom}}

\title{Multiple equivalent characterizations of the Schrödinger equation}
\author{Gabriele Carcassi}
\category{Quantum mechanics}
\tags{Unitary evolution, Schrödinger equation, probability, entropy}

\begin{document}

\maketitle

\begin{abstract}
	We show that the Schrödinger equation is equivalent to determinism and reversibility, which can be characterized in multiple ways: unitary evolution, one-to-one correspondence between past and future states, preservation of conditional probability or conservation of information entropy. The result holds in the infinite-dimensional case.
\end{abstract}

\section{Introduction}

In quantum mechanics, time evolution for an isolated system is postulated to follow the Schrödinger equation. This is equivalent to saying that evolution is represented by a strongly continuous one-parameter group of unitary operators. Here we show that the same evolution can be characterized in more physical terms as deterministic and reversible evolution: the future can be predicted from the past and the past reconstructed from the future. This same assumption can be implemented by requiring a one-to-one correspondence between states, preservation of conditional probability or conservation of information entropy.

\section{Mathematical prerequisites}

We first show that an affine continuous map is a bijection on pure states if and only if it is a bijection on the whole ensemble space, and if and only if it is implemented by either a unitary or an anti-unitary.

\begin{prop}[Affine pure-state correspondence]\label{AFF-PURE}
	Let $D(\mathcal{H})$ be the space of density operators of a Hilbert space $\mathcal{H}$. Let $T:D(\mathcal{H})\to D(\mathcal{H})$ be a trace-norm continuous affine map. Then the following are equivalent:
	\begin{itemize}
		\item $T$ is bijective on the pure states;
		\item $T$ is bijective on all density operators;
		\item there exists either a unitary or an anti-unitary operator $W:\mathcal{H}\to\mathcal{H}$ such that $T(\rho)=W\rho W^\dagger$ for every $\rho\in D(\mathcal{H})$.
	\end{itemize}
\end{prop}

\begin{proof}
	Suppose first that $T$ maps pure states bijectively onto pure states. The affine map on density operators extends to a positive trace-preserving real-linear map on the self-adjoint trace-class operators. For a positive trace-class operator $A\neq 0$, set
	\begin{equation}
		\widetilde{T}(A)=\tr(A)T\left(\frac{A}{\tr(A)}\right).
	\end{equation}
	Affinity makes this extension additive on the positive cone, and every self-adjoint trace-class operator is the difference of two positive ones. This gives a continuous real-linear map on the self-adjoint trace-class operators.
	
	Pure states are exactly the positive trace-one rank-one operators. Therefore $\widetilde{T}$ maps rank-one positive operators bijectively onto rank-one positive operators. The standard linear rank-one preserver result gives two non-degenerate possibilities: $\widetilde{T}(A)=WAW^\dagger$ or, after choosing an orthonormal basis, $\widetilde{T}(A)=WA^{\mathsf T}W^\dagger$. The remaining possibility would map all pure states to the same pure state and is excluded by bijectivity. Positivity and preservation of the trace force $W$ to preserve the norm and the bijection on pure states forces its range to be the whole Hilbert space. Consequently, $W$ is unitary in the first case. In the second case, the transpose can be absorbed into complex conjugation and the resulting operator is anti-unitary. Thus the first condition implies the third.

	A unitary or anti-unitary operator is invertible, and conjugation by it maps density operators onto density operators. Therefore the third condition implies the second.

	Lastly, suppose $T$ maps the ensemble space bijectively onto itself. Since $T$ is affine and bijective, its inverse is affine as well. Affine bijections map extreme points bijectively onto extreme points, and the extreme points of $D(\mathcal{H})$ are exactly the pure states. Therefore the second condition implies the first, completing the proof.
\end{proof}

\begin{remark}
	In finite dimensions, it suffices that $T$ is injective on the pure states $P(\mathcal{H})$, as the compactness of $P(\mathcal{H})$ will force $T$ to be bijective on $P(\mathcal{H})$. In infinite dimensions, this does not work: the map $|n\> \mapsto |n+1\>$, which shifts the energy levels of a harmonic oscillator, is a counterexample.
\end{remark}

\section{Conditions}

The following condition represents the standard structure of quantum states on Hilbert space.
\begin{condition}[Quantum states]{QST:H}
	The state space of a quantum system is represented by the projective space $\mathcal{P}(\mathcal{H})$ of a separable complex Hilbert space $\mathcal{H}$. The ensemble space is represented by the density operators $\mathcal{D}(\mathcal{H})$. The conditional probability is given by the Born rule $p(\psi|\phi) = \frac{\< \phi | \psi \>\< \psi | \phi \>}{\< \psi | \psi \>\< \phi | \phi \>}$. The entropy is given by the von Neumann entropy: $S(\rho) = -\tr(\rho \log \rho)$.
\end{condition}

Time evolution must preserve statistical mixtures.
\begin{condition}[Evolution of mixtures]{[B]EVMIX}
	The evolution of a statistical mixture is the statistical mixture of the evolutions. That is, the evolution from each $t_0$ to each $t_1$, with $t_1\geq t_0$, is given by an affine map $T_{t_0,t_1}:D(\mathcal{H})\to D(\mathcal{H})$ such that $T_{t_0,t_0}=I$, $T_{t_1,t_2}\circ T_{t_0,t_1}=T_{t_0,t_2}$ and $(\rho,t_0,t_1)\mapsto T_{t_0,t_1}(\rho)$ is jointly continuous.
\end{condition}

We are going to restrict the dynamics to the time-independent case.
\begin{condition}[Time-independent evolution]{EVTI}
	Evolution depends only on elapsed time. That is, $T_{t_0,t_1}=T_{t_1-t_0}$.
\end{condition}

Lastly, these are the five conditions whose equivalence we want to prove.
\begin{condition}[Schrödinger equation]{DR-SCEQ}
	The evolution follows the equation
	\begin{equation}
		\imath\hbar\frac{d}{dt}\psi(t)=H\psi(t),
	\end{equation}
	where $H$ is a self-adjoint operator and $\psi(0)\in\Dom(H)$.
\end{condition}

\begin{condition}[Unitary evolution]{DR-UNIT}
	Time evolution is represented by a strongly continuous one-parameter group of unitary operators. That is, $T_t(\rho)=U_t\rho U_t^\dagger$, where $U_t:\mathcal{H}\to\mathcal{H}$, $U_0=I$, $U_{t+s}=U_tU_s$ and $U_t^\dagger U_t=U_tU_t^\dagger=I$.
\end{condition}

\begin{condition}[Bijective map]{DR-EV}
	Distinct past states map to distinct future states and vice-versa. That is, if $\psi_I$ and $\phi_I$ are two initial states, $\psi_F$ and $\phi_F$ are the corresponding evolved future states and $\psi_P$ and $\phi_P$ are the corresponding reconstructed past states, then $\psi_I = \phi_I \iff \psi_F = \phi_F \iff \psi_P = \phi_P$.
\end{condition}

\begin{condition}[Preservation of conditional probability]{DR-PROB}
	Conditional probability remains the same across past and future states. That is, if $\psi_I$ and $\phi_I$ are two initial states, $\psi_F$ and $\phi_F$ are the corresponding evolved future states and $\psi_P$ and $\phi_P$ are the corresponding reconstructed past states, then $p(\psi_I \mid \phi_I) = p(\psi_F \mid \phi_F) = p(\psi_P \mid \phi_P)$.
\end{condition}

\begin{condition}[Conservation of information entropy]{DR-INFO}
	Future and past ensembles have the same information entropy. That is, if $\rho_I$ is an initial ensemble, $\rho_F$ is the corresponding evolved future ensemble and $\rho_P$ is the corresponding reconstructed past ensemble, then $S(\rho_I) = S(\rho_F) = S(\rho_P)$.
\end{condition}

\section{Theorem}

Now we show the equivalence of the time evolution characterizations.

\begin{thrm}
	Over \ref{QST:H}+\ref{[B]EVMIX}+\ref{EVTI}, conditions \ref{DR-SCEQ}, \ref{DR-UNIT}, \ref{DR-EV}, \ref{DR-PROB} and \ref{DR-INFO} are equivalent.
\end{thrm}

\begin{proof}
	We first show that each physical condition gives unitary evolution.

	\ref{DR-EV} $\implies$ \ref{DR-UNIT}. By \ref{DR-EV}, $T_t$ maps pure states bijectively onto pure states. By \ref{[B]EVMIX}, it is continuous and affine. Proposition \ref{AFF-PURE} shows that each $T_t$ is implemented by either a unitary or an anti-unitary operator. The anti-unitary case is excluded by composition. In fact, $T_t=T_{t/2}\circ T_{t/2}$, and the composition of two maps implemented by an anti-unitary is implemented by a unitary. Therefore every $T_t$ is implemented by a unitary.

	The composition law gives a continuous projective unitary representation of elapsed time. The phases can be chosen so that $U_0=I$ and $U_{t+s}=U_tU_s$. Joint continuity of $T_t$ gives strong continuity of $U_t$. The inverses define $U_{-t}=U_t^\dagger$, yielding a strongly continuous one-parameter unitary group. Therefore \ref{DR-UNIT} holds.

	\ref{DR-PROB} $\implies$ \ref{DR-EV}. Two states are the same if and only if their conditional probability is one. Therefore
	\begin{equation}
		\psi_I=\phi_I
		\iff p(\psi_I\mid\phi_I)=1
		\iff p(\psi_F\mid\phi_F)=1
		\iff \psi_F=\phi_F,
	\end{equation}
	and the same argument applies to the reconstructed past states. Therefore \ref{DR-PROB} implies \ref{DR-EV}, which implies \ref{DR-UNIT} by the previous argument.

	\ref{DR-INFO} $\implies$ \ref{DR-UNIT}. Pure ensembles are exactly the ensembles with zero von Neumann entropy. Since the entropy is the same in both directions, pure ensembles map bijectively onto pure ensembles. Affinity and continuity then allow us to apply Proposition \ref{AFF-PURE}, and the same composition argument gives \ref{DR-UNIT}.

	We now show that unitary evolution gives all three physical conditions. Given the evolved future states $\psi_F=U_t\psi_I$ and $\phi_F=U_t\phi_I$, the corresponding reconstructed past states are $\psi_P=U_t^\dagger\psi_F=\psi_I$ and $\phi_P=U_t^\dagger\phi_F=\phi_I$. Since $U_t$ and $U_t^\dagger$ are bijective, \ref{DR-UNIT} implies \ref{DR-EV}. Moreover,
	\begin{equation}
		p(\psi_I\mid\phi_I)=p(\psi_F\mid\phi_F)=p(\psi_P\mid\phi_P),
	\end{equation}
	so \ref{DR-PROB} holds. Lastly, let $\rho_F=U_t\rho_IU_t^\dagger$ and $\rho_P=U_t^\dagger\rho_FU_t=\rho_I$. Unitary conjugation preserves eigenvalues, including multiplicities. Since the von Neumann entropy depends only on those eigenvalues,
	\begin{equation}
		S(\rho_I)=S(\rho_F)=S(\rho_P),
	\end{equation}
	and \ref{DR-INFO} holds.

	It remains to connect the unitary group to the Schrödinger equation. By Stone's theorem, every strongly continuous one-parameter unitary group has a unique self-adjoint generator $H$ such that
	\begin{equation}
		U_t=e^{-\imath tH/\hbar}.
	\end{equation}
	For every $\psi(0)\in\Dom(H)$, the curve $\psi(t)=U_t\psi(0)$ is differentiable and satisfies \ref{DR-SCEQ}. Conversely, a time-independent self-adjoint Hamiltonian defines through the same expression a strongly continuous one-parameter unitary group. Therefore \ref{DR-SCEQ} and \ref{DR-UNIT} are equivalent, completing the proof.
\end{proof}


\section{Additional remarks}

\textbf{The Hamiltonian domain.} For an unbounded self-adjoint Hamiltonian, the Schrödinger equation is not defined on every vector of the Hilbert space. The unitary operators $e^{-\imath tH/\hbar}$ are defined everywhere, but the derivative is guaranteed only for initial vectors in $\Dom(H)$. The group formulation is therefore the more general statement of time evolution.

\textbf{The phase of the unitary.} Evolution of density operators does not determine the overall phase of $U_t$. Consequently, its self-adjoint generator is determined only up to an additive multiple of the identity. This does not change states, conditional probabilities or entropy.

\section{Conclusion}

We have seen that the Schrödinger equation can be characterized in different ways, based on bijective maps of pure states, preservation of probability and conservation of entropy. Conceptually, all these are equivalent to saying that the system is isolated, as future and past states are determined, and only determined, by the present state of the system.

\end{document}
